\section{Introduction}
Alzheimer's disease (AD) and type 2 diabetes mellitus (T2DM) represent two of the most pressing chronic health challenges we face today \cite{vb2025global, lopez2024recent}, each carrying its own significant public health burden while also sharing surprising connections with one another \cite{a2014linking}. Today, AD stands as the primary cause of dementia in older adults \cite{alzheimer20152015, mezey2000advance}. More than 55 million people worldwide are living with dementia, and experts predict this could surge to somewhere between 139 and 150 million by 2050 \cite{schwarzinger2022forecasting, logroscino2020prevention}. T2DM shows a similarly alarming trend, with roughly 589 million adults as of 2024 (that's about 11.1\% of the adult population) living with the condition, and projections suggest we could see around 853 million cases by mid-century \cite{kalyani2025diagnosis}. What is particularly striking is that over 90\% of all diabetes cases are type 2 \cite{kalyani2025diagnosis, zheng2018global}, largely driven by our ageing populations and modern lifestyle factors. As life expectancy continues to climb, there are statistically more people dealing with both metabolic and neurodegenerative diseases simultaneously \cite{procaccini2016role}, which makes understanding the relationship between AD and T2DM increasingly urgent for biomedical AI development.

 Epidemiological research has consistently shown that having diabetes substantially increases the risk of cognitive decline and dementia, including AD \cite{barbagallo2014type}. A recent 2024 meta-analysis showed that diabetic patients face about 59\% higher dementia risk compared to people without diabetes\cite{barbagallo2014type, bellia2022diabetes}, which reinforces what earlier studies had found that T2DM increases an individual's risk of cognitive disorders by roughly 1.3 to 1.9 times \cite{kciuk2024alzheimer}. When we dig into the molecular details, the picture gets even more interesting. The chronic high blood sugar and insulin resistance that defines T2DM actually mirrors some of the same pathophysiological processes we see in AD \cite{blazquez2014insulin}. This similarity has led some researchers to provocatively label AD as ``type 3 diabetes" though this remains controversial \cite{abdelgadir2017effect, de2014type, kroner2009relationship, leszek2017type}. The idea emphasizes the overlapping features: impaired insulin signaling in the brain, chronic inflammation, and oxidative stress \cite{kciuk2024alzheimer}. Both conditions also share common risk factors such as midlife obesity and high blood pressure \cite{kciuk2024alzheimer}.

Despite the wealth of published research on AD–T2DM connections, Large Language Models (LLMs) on their own have real trouble delivering reliable medical answers \cite{lacerdaevaluation}. Today's state-of-the-art (SOTA) LLMs can certainly generate impressively fluent responses, but they frequently hallucinate facts or fabricate citations that do not exist, which is a serious safety concern in healthcare settings \cite{wang2025trustworthy}. Even powerful models like GPT-4 have inherent knowledge limitations that mean they might miss crucial domain-specific details or misrepresent the latest findings. This is where Retrieval-Augmented Generation (RAG) has become the go-to solution, helping ground LLM outputs in external, verified knowledge \cite{wu2025multirag}. By pulling in relevant documents or facts when answering a query, RAG significantly improves factual accuracy \cite{wu2025multirag}. However, conventional RAG systems that retrieve free-text passages can still drag in irrelevant or contradictory information, which the LLM might then mistakenly weave into its answer \cite{wu2025multirag}. This becomes especially problematic for knowledge-intensive questions that require multiple reasoning steps, where models need to piece together complex biomedical relationships. 

\begin{table}[ht]
\centering
\caption{Examples of triples used for the knowledge graphs}
\scriptsize
\setlength{\tabcolsep}{4pt}
\renewcommand{\arraystretch}{1.1}
\resizebox{\linewidth}{!}{%
\begin{tabular}{p{.16\linewidth} p{.84\linewidth}}
\hline
\textbf{KG} & \textbf{Example} \\
\hline
\textbf{$\mathbb{G}_1$} & [{"Entity 1": "T2DM", "Relationship": "was associated with", "Entity 2": "decreased forced expiratory volume in 1s (FEV1)"}] \\
\textbf{$\mathbb{G}_2$} & [{"Entity 1": "Alzheimer's disease CSF", "Relationship": "is associated with", "Entity 2": "neuroinflammation"}] \\
\textbf{$\mathbb{G}_3$} & [{"Entity 1": “Insulin-like growth factor ", "Relationship": "influences", "Entity 2": “cognitive functions"}] \\
\hline
\end{tabular}}
\label{tab:kg_examples}
\end{table}

Structured knowledge bases offer a more promising alternative. Knowledge graphs (KGs) organize facts as subject–relation–object triples $(E_1, R, E_2)$ as seen in Table~\ref{tab:kg_examples}, where $E_1$ is the subject entity, $E_2$ is the object entity, and $R$ represents their relationship. This structured representation provides clear semantics and traceable origins, making it easier to verify how an answer was constructed. Yet, critical questions remain about how best to leverage KGs in healthcare AI: Does curating domain-specific KGs from biomedical literature actually improve answer quality in RAG-enhanced LLMs? Should we build narrow, disease-focused graphs or broader cross-domain ones? And how do these design choices interact with model parameters like decoding temperature?

In this research, we investigate the role of domain-specific KGs curated from PubMed abstracts in RAG-enhanced healthcare LLMs. We ask whether such KGs improve the factual accuracy and evidential support of answers compared to no-RAG baselines (\textbf{RQ1}). We further examine how curation scope affects question answering: do focused, disease-specific KGs (G1 for T2DM, G2 for AD) outperform a broader merged KG (G3) on single-hop and multi-hop probes, or vice versa (\textbf{RQ2})? We also investigate the robustness of RAG gains across different decoding temperatures (0, 0.2, 0.5) and test for statistical significance of improvements (\textbf{RQ3}). Finally, we explore whether we can enhancing the components of the CoDe-KG pipeline to yield higher-quality graphs that translate into better downstream QA performance (\textbf{RQ4}).


% 
% We also look into how the type of curated KG affects the output of the LLM.  We developed a RAG-enhanced LLM that uses a custom-built KG as external knowledge to answer questions at the intersection of AD and T2DM. We constructed our knowledge base from three sets of triples: $\mathbb{G}_1$ (abstract-derived triples targeting T2DM), $\mathbb{G}_2$ (abstract-derived triples targeting AD), and $\mathbb{G}_3$ (abstract-derived triples spanning both T2DM and AD). To identify the most relevant cross-domain knowledge, we used a cosine similarity threshold of 0.65 to find AD and T2DM triples with meaningful overlap, creating an intersection set we refer to as $\mathbb{G}_1$ ∩ $\mathbb{G}_2$. Both $\mathbb{G}_1$ ∩ $\mathbb{G}_2$ and $\mathbb{G}_3$ serve as the foundation for our knowledge graphs, allowing us to compare focused versus comprehensive knowledge representations.

\subsection{Contributions}
In this research, we investigate the role of domain-specific KGs curated from abstracts on PubMed for RAG-enhanced healthcare LLMs. 
\begin{itemize}
    \item We build three abstract-derived KGs with a common schema: \textbf{$\mathbb{G}_1$} (T2DM-focused), \textbf{$\mathbb{G}_2$} (AD-focused), and \textbf{$\mathbb{G}_3$} (combined AD+T2DM).
   \item We introduce probe sets spanning single-hop clinical facts (e.g., drug$\rightarrow$outcome, gene$\rightarrow$disease) and multi-hop mechanisms (e.g., insulin signaling$\rightarrow$neuroinflammation$\rightarrow$cognition)
   \item We evaluate 7 LLMs \emph{with} and \emph{without} KG-RAG across varying temperatures $0$, $0.2$, and $0.5$, measuring F$_1$score, and accuracy across the probes, and we test for significance of RAG gains and of curation scope ($\mathbb{G}_1$/$\mathbb{G}_2$ vs.\ $\mathbb{G}_3$)
   \item We improve the co-reference of the CoDe-KG pipeline.
\end{itemize}









% \section{Introduction}
% Knowledge graphs (KG) are systems that .... They are of many timees including .... The common appearance of the KG includes the Entity1, Relationship Entity2. The importance of this in healthcare cannot be over emphasised.

% Recent studies have looked into Alzheimer's disease (AD) and its relation to T2DM diabetes.  (write on it AD in 2 lines on how it a menace and then also on T2DM and also a menance and how metformin treatments and apoe e4 allele has been typically found in cases)

% Then motivate the need to build a knowledge graph based on these two, to understand the relationships. Then also building trust worthy LLMs that can answer questions on both T2DM and AD which can help us formalize research.

% We plan to test if the RAG systems actually have an impact on the LLMs on a significant level accoring to other research especially in this particular domain and then is there a significance in the kind of knowledge that is being curated. For this we create three graphs $\mathbb{G}_1$ - Graph built on abstracts targeting  T2DM, $\mathbb{G}_2$ - Graph built on Alzhiemers disease and $\mathbb{G}_3$ - Graph built on abstracts targeting both diseases. Based on these graphs, we formulate probes that are used to test LLMs and we judge if the type of knowlegde curation of these graphs has an effect on the way the graph is strcuctured